% The hyperref package is used for handling cross-referencing in the document. It makes links clickable, and can also be used to create bookmarks within the PDF document.
\usepackage{hyperref}

% The graphicx package is used for including graphics files within a LaTeX document. It supports various options that allow you to manipulate the image's dimensions, rotation, and placement. For example, you can use the \includegraphics command to insert an image.
\usepackage{graphicx}

% The subcaption package is used for creating sub-figures or sub-tables within a single float. This allows you to have, for example, multiple images under a single figure caption, each with its own subcaption.
\usepackage{subcaption}

% The tabularx package allows you to create tables with specified column widths. This is especially useful when the content of the table cells needs to be wrapped. The X column type is introduced by this package for this purpose.
\usepackage{tabularx}

% The acro package is used for managing acronyms and abbreviations in your LaTeX document. It allows you to define acronyms once and use them throughout the text while automatically handling their plural forms and capitalization.
\usepackage{acro}

% The fancyhdr package allows you to customize the headers and footers of your LaTeX document. You can add custom text, page numbers, and even images to the header and footer regions.
\usepackage{fancyhdr}

% The lipsum package is used to generate filler text, commonly known as "Lorem Ipsum". This is useful for testing and developing layouts before the final text is available.
\usepackage{lipsum}




